\chapter{A Concrete Naming Certificate for $\PseudoArcUp$}
\label{ch:concrete-instances-pseudo-arc-namecert}
\origin{human}

Following Bing's pseudo-arc construction, the $\PseudoArcUp$ packet records a hereditarily indecomposable chainable compact metric continuum through finite inverse-limit rows rather than through a host continuum object. It sits over the compact-metric source of \autoref{ch:concrete-instances-compactmetric-namecert}, the finite projection discipline of \autoref{ch:concrete-instances-inverse-limit-metric-namecert}, the complete-metric boundary of \autoref{ch:concrete-instances-completemetric-namecert}, the stream-window surface of \autoref{ch:concrete-instances-streamname-namecert}, and the regular rational readback of \autoref{ch:concrete-instances-regseqrat-namecert}. The packet names only displayed finite chain covers, projection refinements, and indecomposability obstruction rows before any continuum-facing consumer may read it.

\begin{definition}[Pseudo-arc carrier]
\label{def:pseudo-arc-carrier}
A $\PseudoArcUp$ carrier is a finite $\BHist$ packet
$$
P=(K,I,C,D,R,E,H,T,G,N)
$$
with the following displayed rows:
\begin{enumerate}
\item $K$ is the $\CompactMetricUp$ source row, exposing finite metric covers, compact-metric provenance, and distance readback for the current continuum window.
\item $I$ is the $\InverseLimitMetricUp$ row, recording a finite prefix of compact metric continua, bonding-map windows, and projection-compatibility checks.
\item $C$ is the chainability row: for each displayed scale it lists a finite chain cover and its adjacency ledger over the same compact metric source.
\item $D$ is the hereditary-indecomposability row, recording finite obstruction tests for splitting a displayed subcontinuum window into two proper subcontinua.
\item $R$ is the $\RegSeqRatUp$ readback row for the diameters, overlap tolerances, and bonding-map moduli inspected by $C$ and $D$.
\item $E$ is the terminal $\RealUp$ seal for the requested tolerance budget.
\item $H$ records componentwise $\hsame$ transport for compact-metric, inverse-limit, chainability, indecomposability, readback, seal, provenance, and naming rows.
\item $T$ records displayed $\Cont$ replay from compact-metric cover requests through inverse-limit projections, chain covers, obstruction tests, and regular rational readback.
\item $G$ is the $\Pkg$ provenance over $\PseudoArcUp$, $\CompactMetricUp$, $\InverseLimitMetricUp$, $\CompleteMetricUp$, $\StreamNameUp$, $\RegSeqRatUp$, $\RealUp$, $\BHist$, $\hsame$, $\Cont$, and $\NameCert$.
\item $N$ is the local $\NameCert_{\PseudoArcUp}$ row.
\end{enumerate}
The classifier compares accepted packets only through $K,I,C,D,R,E,H,T,G,N$. It has no coordinate for a host pseudo-arc, quotient continuum equality, a choice-selected inverse-limit point, open-cover compactness, or an ambient completeness theorem.
\end{definition}

\begin{theorem}[Pseudo-arc NameCert obligations]
\label{thm:pseudo-arc-namecert-obligations}
For every accepted $\PseudoArcUp$ carrier $P=(K,I,C,D,R,E,H,T,G,N)$, the local $\NameCert_{\PseudoArcUp}$ surface consists exactly of compact-metric source admission through $K$, finite inverse-limit prefix admission through $I$, displayed chain-cover coverage through $C$, hereditary-indecomposability obstruction exposure through $D$, regular rational tolerance readback through $R$, terminal real sealing through $E$, componentwise transport through $H$, displayed replay through $T$, provenance through $G$, and local naming through $N$.

Every continuum-facing consumer must read these rows before it may use the packet as pseudo-arc data, and no obligation row is discharged by a host pseudo-arc theorem, quotient equality of continua, selected inverse-limit representatives, open-cover compactness, countable choice, or ambient $\RealUp$ completeness.
\end{theorem}
\begin{proof}
Project the carrier of \autoref{def:pseudo-arc-carrier}. The compact source $K$ and inverse-limit prefix $I$ are the only continuum-source rows, so every consumer begins with displayed compact-metric covers and finite projection compatibility.

The chainability and hereditary-indecomposability rows are finite. The row $C$ lists chain covers at inspected scales, while $D$ lists obstruction windows preventing a displayed subcontinuum split. The regular rational row $R$ supplies the tolerance readback needed by both rows before the real seal $E$ is consumed.

Transport and replay keep the packet local. The rows $H$ and $T$ preserve the same compact-metric, inverse-limit, chainability, obstruction, readback, and seal data under $\hsame$ and $\Cont$, while $G$ and $N$ provide provenance and naming.

Since the carrier has no coordinate for a host pseudo-arc, quotient continuum equality, selected inverse-limit point, open-cover compactness, countable choice, or ambient real completeness, the local NameCert surface is exhausted by the displayed rows.
\end{proof}
\leanstmt{BEDC.Derived.PseudoArcUp}

\falsifiablePrediction{Every accepted $\PseudoArcUp$ read exposes compact-metric cover rows, finite inverse-limit prefix rows, chain-cover rows, hereditary-indecomposability obstruction rows, regular rational tolerance readback, a Real seal, componentwise $\hsame$ transport, displayed $\Cont$ replay, $\Pkg$ provenance, and local $\NameCert$ before a continuum-facing consumer can inspect the packet.}

\independenceWitness{compactmetric, inverse-limit-metric, completemetric, streamname, regseqrat: $\PseudoArcUp$ is not bijective to any listed sibling because this packet binds chainability rows, hereditary-indecomposability obstruction rows, finite inverse-limit prefixes, compact-metric source rows, readback, transport, replay, provenance, and local naming in one carrier.}

\closureat{PseudoArcUp}{seedStr}

\begin{closurestatus}{\PseudoArcUp}
  \constructivestory{A $\PseudoArcUp$ packet is generated from CompactMetricUp source rows, finite InverseLimitMetricUp prefix rows, chain-cover rows, hereditary-indecomposability obstruction rows, RegSeqRatUp tolerance readback, a RealUp seal, componentwise hsame transport, displayed Cont replay, Pkg provenance, and a local NameCert row.}
  \theoryclosure{\seedClosure}
  \scopeclosed{\autoref{def:pseudo-arc-carrier} and \autoref{thm:pseudo-arc-namecert-obligations} seed the finite pseudo-arc boundary over compact-metric, inverse-limit, chainability, obstruction, readback, and naming rows.}
  \formalstatus{\formalTargetV}
  \bridgestatus{none}
  \notclaimed{This chapter does not claim a host pseudo-arc theorem, quotient continuum equality, selected inverse-limit representatives, open-cover compactness, countable choice, Lean-checked obligations, or a standard bridge.}
  \upgradepath{Obligation closure requires checked carrier admission, finite inverse-limit prefix admission, chain-cover coverage, hereditary-indecomposability obstruction exposure, tolerance readback, transport, replay, provenance, and local naming for BEDC.Derived.PseudoArcUp.}
\end{closurestatus}
